
%-------------------------------------------------------------------------------
\section{Real-World Attacks} \label{sec:real_world}
%-------------------------------------------------------------------------------

In this section, we present a real-world attack case that was identified and successfully carried out, inspired by our validation results.
The vulnerability exploited by this attack is widely present on the Android platform, affecting a large number of mobile applications that adopt WeChat Pay and Alipay as third-party payment methods, and has been officially confirmed by the security teams of both WeChat Pay and Alipay.

\subsection{A Vulnerability in Real-World}

Our formal analysis reveals that variant \Rom{2} of the order tampering attack demonstrates that an adversary can not only tamper with orders over the communication channel between \emph{MF} and the \emph{MS}, but also conduct similar attacks on the communication channel between the \emph{MF} and the \emph{PS-App}.
Based on the above findings, we identify security implementation flaws in some merchant applications when integrating WeChat Pay or Alipay.

When initiating a payment request, a merchant application needs to transmit order parameters to the payment service application via the Android Intent mechanism.
As a core inter-component communication mechanism provided by the Android system, Intent serves as the primary means by which the WeChat Pay and Alipay receive payment parameters from merchant applications.
A merchant application typically constructs an Intent containing order information and uses it to launch the corresponding payment application.

It should be noted that Intents can be classified into two categories: explicit Intents and implicit Intents. 
An explicit Intent specifies the fully qualified package name of the target application, whereas an implicit Intent only declares the action, category, and data, allowing the system to match and select a responding component based on Intent Filters.
The security risk primarily arises from the misuse of implicit Intents.

When invoking payment functionality, a merchant application should use an explicit Intent to explicitly designate the official WeChat Pay or Alipay as the target, thereby ensuring that the payment request is correctly routed to a trusted payment application.
However, we find that some merchant applications use implicit Intents in practice.
This enables an attacker to hijack payment requests by installing a malicious application on the victim's device that declares the same Intent-handling capabilities.


\subsection{Order Tampering Attack}

\begin{figure}[t]
  \centering
  \includegraphics[width=\linewidth]{img/black.png}
  \caption{Real-World Attack Setup}
  \label{fig:real_world_attack_devices}
\end{figure}

By exploiting the aforementioned vulnerability, an attacker can carry out the variant~\Rom{2} of the order tampering attack.
To carry out this attack in the real world, we developed two Android applications.
One is a malicious application~(Malicious App), which is used to intercept and tamper with payment requests.
The other is an attacker application (Attacker App), which is responsible for generating payment orders for the attacker.
We use two Android devices, as illustrated in Fig.~\ref{fig:real_world_attack_devices}: one serves as the victim device (Victim Device), and the other serves as the attacker device (Attacker Device).
On the victim device, we install the malicious application, a vulnerable merchant application, and the payment app.
The attacker application is installed on the attacker device.

\begin{figure}[t]
  \centering
  \includegraphics[width=\linewidth]{img/black.png}
  \caption{Real-World Attack Flow}
  \label{fig:real_world_attack_flow}
\end{figure}

Fig.~\ref{fig:real_world_attack_flow} illustrates the workflow of the real-world order tampering attack.
The attack procedure can be summarized as follows:
\begin{enumerate} 
  \item The user initiates a payment within a benign third-party Android application;
  \item The merchant application transmits order parameters to the payment application using an implicit Intent;
  \item The malicious application induces the user to select itself as the default handler (or automatically sets it via accessibility services);
  \item The malicious application intercepts the Intent and forwards the captured payment parameters to the Attacker App;
  \item The Attacker App constructs a new payment order based on attacker-specified parameters and returns it to the malicious application;
  \item The malicious application embeds the tampered order parameters back into the Intent and launches the payment application;
  \item The payment application receives the tampered order parameters and completes the payment process.
\end{enumerate}
As a result, the attack successfully achieves order tampering, causing the user to ultimately pay for an order specified by the attacker.

\subsection{Impact Assessment}
To assess the real-world impact of the identified vulnerability, we conducted a manual measurement study on Android applications and games that integrate the WeChat Pay or Alipay SDK.
Our dataset consists of 150 popular Android applications and 50 games, covering up to the three most recent versions of each application and the latest version of each game, resulting in a total of 461 app versions.

Among these versions, 419 support WeChat Pay and 401 support Alipay.
Our analysis reveals that 64 WeChat Pay–enabled versions and 72 Alipay-enabled versions are vulnerable.
In total, 116 app versions (more than 25\% of the analysed dataset) are found to be vulnerable, meaning that at least one of the integrated payment SDKs is affected.

The vulnerable applications span a wide range of categories, including e-commerce, social networking, online education, and entertainment.
More importantly, the cumulative number of downloads of these vulnerable applications reaches 265.85 billion, indicating that the vulnerability potentially affects a very large user base and poses a significant security risk in practice.

\subsection{Responsible Disclosure}
In December 2025, we responsibly reported this vulnerability to the security teams of WeChat Pay and Alipay, both of which acknowledged the issue and treated it with high priority.
As of the time of writing, Alipay has committed to fixing the vulnerability by the end of January 2026, and WeChat Pay is likewise actively working on a patch.
In addition, we contacted the developers of xxx affected Android applications and games, recommending the use of explicit Intents to ensure the secure transmission of payment requests.
We received positive responses from xxx developers, who indicated that they would update their applications promptly to address the issue.
